\chapter{Introduction}
\label{ch:introduction}

% Establish the problem and research agenda; reserve technical detail for Chapter 2.
\section{Proteins as Dynamic Ensembles}
% Explain why conformational populations matter for function and disease.
% proteins are dynamic polypeptides → protein folding grand challenge → folding versus function (dynamics) → case studies showing why we need to measure and model populations - where current structural methods fall short.


\section{Measuring Protein Dynamics}
% Introduce HDX-MS as the focal measurement and why its structural interpretation is indirect.
% experimental access to protein dynamics → HDX-MS as a solution-phase probe → peptide-level deuterium uptake over labelling time → dependence on intrinsic exchange chemistry and structural protection → different structural explanations can produce similar uptake → need for candidate ensembles to interpret the measurement.


\section{Estimating Protein Dynamics}
% Introduce molecular dynamics and learned generators, including diffusion models.
% candidate ensembles → MD and its limits → enhanced sampling → learned generators including diffusion → why experimental constraints are needed to assess populations. Keep the equations and detailed algorithms in Chapter 2, as your outline already proposes.

\section{Integrating Experiment and Ensembles}
% Frame the inverse problem and the role of structure-to-observable forward models, with HDX-MS as the case.
% indirect measurements and the inverse problem → candidate ensembles and forward prediction → reweighting → prior and MaxEnt → joint population–parameter fitting → why the resulting models require validation.

\section{Trusting Integrative Ensemble Models}
% Motivate ambiguity, missing states, model error, overfitting, and interpretability.
% non-uniqueness → candidate coverage and decoys → model–population compensation → generalisation and regularisation → uncertainty, interpetation and testing of structural hypotheses

\section{Research Questions and Objectives}
% State the questions that connect validation, joint fitting, and guided generation.

\section{Thesis Outline and Contributions}
% Map each research chapter to its question and contribution; distinguish ongoing work.
